\documentclass[a4paper,12pt,titlepage,finall]{article}

\usepackage[T1,T2A]{fontenc}     % форматы шрифтов
\usepackage[utf8x]{inputenc}     % кодировка символов, используемая в данном файле
\usepackage[russian]{babel}      % пакет русификации
\usepackage{tikz}                % для создания иллюстраций
\usepackage{pgfplots}            % для вывода графиков функций
\usepackage{geometry}		 % для настройки размера полей
\usepackage{indentfirst}         % для отступа в первом абзаце секции
\usepackage{minted}
\usepackage{amsmath}

% выбираем размер листа А4, все поля ставим по 3см
\geometry{a4paper,left=30mm,top=30mm,bottom=30mm,right=30mm}

\setcounter{secnumdepth}{0}      % отключаем нумерацию секций

\usepgfplotslibrary{fillbetween} % для изображения областей на графиках

\begin{document}
% Титульный лист
\begin{titlepage}
    \begin{center}
	{\small \sc Московский государственный университет \\имени М.~В.~Ломоносова\\
	Факультет вычислительной математики и кибернетики\\}
	\vfill
	{\Large \sc Отчет по заданию №6}\\
	~\\
	{\large \bf <<Сборка многомодульных программ. \\
	Вычисление корней уравнений и определенных интегралов.>>}\\ 
	~\\
	{\large \bf Вариант 11 / 5 / 4}
    \end{center}
    \begin{flushright}
	\vfill {Выполнил:\\
	студент 102 группы\\
	Соловьев~А.~И.\\
	~\\
	Преподаватели:\\
	Калинин~Г.~М.,\\Цыбров~Е.~Г.}
    \end{flushright}
    \begin{center}
	\vfill
	{\small Москва\\2024}
    \end{center}
\end{titlepage}

% Автоматически генерируем оглавление на отдельной странице
\tableofcontents
\newpage

\section{Постановка задачи}

Необходимо найти площадь фигуры, ограниченной данными тремя кривыми, лежащей на заданном отрезке. Гарантируется, что кривые попарно пересекаются на данном отрезке. 
\begin{itemize}
    \item Для нахождения точек пересечений требуется реализовать численный метод нахождения корней. Решение реализует два метода (на выбор): метод Ньютона (касательных) и метод хорд.
    \item Для нахождения площади фигуры требуется реализовать метод численного интегрирования. Решение реализует два метода (на выбор): метод трапеций и метод Симпсона (парабол).
\item Отрезок для применения метода нахождения корней должен быть вычислен аналитически.
\end{itemize}
Функции, задающие кривые, и отрезок, на котором лежит фигура, задаются перед сборкой в файле. Функции задаются в обратной польской записи.\\\\
Пример:
\begin{minted}[frame=single, linenos, breaklines=true, fontsize=\small]{text}
0.0 4.0
2 x 4 / tan -
x
0.2 pi *
\end{minted}

\newpage

\section{Математическое обоснование}

В данном разделе проводится анализ заданного набора кривых, приводятся их графики (рис.~\ref{plot1}),
обоснование выбора значений $\varepsilon_1$ и $\varepsilon_2$, а также отрезков для поиска точек
пересечения кривых.

\subsection{Минимальные условия сходимости}
\begin{itemize}
    \item $\exists \varepsilon: f(\varepsilon) = 0, f'(\varepsilon) \neq 0$
    \item $f(x)$ непрерывно дифференцируема на отрезке
    \item начальная точка $x_0$ достаточна близка к $\varepsilon$
\end{itemize}

Для метода Ньютона желательно соблюсти все условия Фурье:
\begin{itemize}
    \item $f'(x) \neq 0$ на отрезке
    \item $f''(x)$ не меняет знак на отрезке
    \item начальная точка $x_0$ с того конца отрезка, где $f(x_0)\cdot f''(x_0) > 0$,
\end{itemize}
так как при соблюдении лишь минимальных условий метод может уйти за пределы отрезка, сойтись к другому корню или не сойтись вовсе.~\cite{math}\\\\

Программа тестировалась на таких наборах данных:\\
1.
\begin{minted}[frame=single, linenos, breaklines=true, fontsize=\small]{text}
0.0 1.0
x
1 x -
x 0.5 - 2 ^
\end{minted}
Аналитическое решение:
$$
\begin{gathered}
    x^2 - x + 0.25 = x\\D = (-2)^2-4\cdot0.25=3\\x_{1,2}=\frac{2\pm \sqrt{3}}{2}=1 \pm \frac{\sqrt{3}}{2},\; x_1 \in [0,\,1]\\\\
    x^2 - x + 0.25 = 1-x,\,x^2-0.75=0\\D = 4\cdot0.75=3\\x_{1,2}=\frac{\pm \sqrt{3}}{2},\;x_2 \in [0, 1]\\\\
x = 1-x;\; x_3=0.5
\end{gathered}
$$

$$
\begin{gathered}
\int x\,dx = \frac{1}{2} x^2 + C\\
\int -x+1\,dx = -\frac{1}{2}x^2+x + C\\
\int (x-0.5)^2\,dx = \frac{1}{3}(x-0.5)^3 + C\\\\
\end{gathered}
$$

$$
\begin{gathered}
\frac{1}{2}\cdot(0.25 - (1 - \frac{\sqrt{3}}{2})^2) \approx 0.116\\
-\frac{1}{2}\cdot(\frac{3}{4} - 0.25) + (\frac{\sqrt{3}}{2} - 0.5) \approx 0.366 - 0.25 = 0.116\\
\frac{1}{3}\cdot (0.5 - \frac{\sqrt{3}}{2})^3 \approx 0.016\\
\frac{1}{3}\cdot (\frac{\sqrt{3}}{2}-0.5)^3 \approx 0.016\\\\
S \approx 0.116 \cdot 2 - 0.016 \cdot 2 = 0.2
\end{gathered}
$$
\\\\

2.
\begin{minted}[frame=single, linenos, breaklines=true, fontsize=\small]{text}
0.0 1.0
x tan
1 x tan -
x tan 10 /
\end{minted}
Аналитическое решение:
$$
\begin{gathered}
tan(x) = 1 - tan(x);\;tan(x) = 0.5;\; x_1 = arctan(0.5);\\
tan(x) = \frac{tan(x)}{10};\;tan(x) = 0; x_2 = 0;\\
1 - tan(x) = \frac{tan(x)}{10};\;tan(x)=\frac{10}{11}; x_3 = arctan(\frac{10}{11})\\
\end{gathered}
$$

$$
\begin{gathered}
\int tan(x)\,dx = -ln|cos\,x| + C\\
\int 1 - tan(x)\,dx = x + ln|cos\;x|+C\\
\int \frac{tan(x)}{10}\,dx = \frac{-ln|cos\,x|}{10} + C\\\\
S = -ln|cos(arctan(0.5)| + (arctan(\frac{10}{11}) + ln|cos(arctan(\frac{10}{11})| - \\- (arctan(0.5)+ln|cos(arctan(0.5|)) + \frac{ln|cos(arctan(\frac{10}{11}))|}{10} = \\
\approx 0.166
\end{gathered}
$$

\newpage

\section{Результаты экспериментов}

В данном разделе приведены результаты проведенных вычислений для двух вышеупомянутых тестовых примеров:
координаты точек пересечения в таблицах и площадь полученной фигуры.

\begin{table}[h]
\centering
\begin{tabular}{|c|c|c|}
\hline
Кривые & $x$ & $y$ \\
\hline
1 и 2 &  0.5 & 0.5 \\
2 и 3 &  0.8660 & 0.1339 \\
1 и 3 & 0.1339 & 0.1339 \\
\hline
\end{tabular}
\caption{Координаты точек пересечения}
\label{table1}
\end{table}


\begin{figure}[h]
\centering
\begin{tikzpicture}
\begin{axis}[% grid=both,                % рисуем координатную сетку (если нужно)
             axis lines=middle,          % рисуем оси координат в привычном для математики месте
             restrict x to domain=-1:2,  % задаем диапазон значений переменной x
             restrict y to domain=-0.25:1.5,  % задаем диапазон значений функции y(x)
             axis equal,                 % требуем соблюдения пропорций по осям x и y
             enlargelimits,              % разрешаем при необходимости увеличивать диапазоны переменных
             legend cell align=left,     % задаем выравнивание в рамке обозначений
             scale=2,                    % задаем масштаб 2:1
             xticklabels={,,},           % убираем нумерацию с оси x
             yticklabels={,,}]           % убираем нумерацию с оси y

% первая функция
% параметр samples отвечает за качество прорисовки
\addplot[green,samples=256,thick,name path=A] {x};
% описание первой функции
    \addlegendentry{$f_1(x)=x$}

% добавим немного пустого места между описанием первой и второй функций
\addlegendimage{empty legend}\addlegendentry{}

% вторая функция
% здесь необходимо дополнительно ограничить диапазон значений переменной x
\addplot[blue,samples=256,thick,name path=B] {1-x};
    \addlegendentry{$f_2(x)=1-x$}

% дополнительное пустое место не требуется, так как формулы имеют небольшой размер по высоте

% третья функция
    \addplot[red,samples=256,thick,name path=C] {(x-0.5)^2};
    \addlegendentry{$f_3(x)=(x-0.5)^2$}

% закрашиваем фигуру
\addplot[blue!20,samples=256] fill between[of=B and C,soft clip={domain=0.5:0.8660}];
\addplot[blue!20,samples=256] fill between[of=A and C,soft clip={domain=0.1339:0.5}];
\addlegendentry{$S=0.1994$}

% Поскольку автоматическое вычисление точек пересечения кривых в TiKZ реализовать сложно,
% будем явно задавать координаты.
\addplot[dashed] coordinates { (0.1339, 0.1339) (0.1339, 0) };
\addplot[color=black] coordinates {(0.1339, 0)} node [label={-135:{\small 0.1339}}]{};

\addplot[dashed] coordinates { (0.5, 0.5) (0.5, 0) };
\addplot[color=black] coordinates {(0.5, 0)} node [label={-90:{\small 0.5}}]{};

\addplot[dashed] coordinates { (0.8660, 0.1339) (0.8660, 0) };
\addplot[color=black] coordinates {(0.8660, 0)} node [label={-45:{\small 0.8660}}]{};

\end{axis}
\end{tikzpicture}
\caption{Плоская фигура, ограниченная графиками заданных уравнений}
\label{plot2}
\end{figure}




















\begin{table}[h]
\centering
\begin{tabular}{|c|c|c|}
\hline
Кривые & $x$ & $y$ \\
\hline
1 и 2 &  0.4636 & 0.5 \\
2 и 3 &  0.7378 & 0.091 \\
1 и 3 & 0 & 0 \\
\hline
\end{tabular}
\caption{Координаты точек пересечения}
\label{table1}
\end{table}

\begin{figure}[h]
\centering
\begin{tikzpicture}
\begin{axis}[% grid=both,                % рисуем координатную сетку (если нужно)
             axis lines=middle,          % рисуем оси координат в привычном для математики месте
             restrict x to domain=-1:2,  % задаем диапазон значений переменной x
             restrict y to domain=-0.25:1.5,  % задаем диапазон значений функции y(x)
             axis equal,                 % требуем соблюдения пропорций по осям x и y
             enlargelimits,              % разрешаем при необходимости увеличивать диапазоны переменных
             legend cell align=left,     % задаем выравнивание в рамке обозначений
             scale=2,                    % задаем масштаб 2:1
             xticklabels={,,},           % убираем нумерацию с оси x
             yticklabels={,,}]           % убираем нумерацию с оси y

% первая функция
% параметр samples отвечает за качество прорисовки
    \addplot[green,samples=256,thick,name path=A] {tan(x * 180 / 3.14159)};
% описание первой функции
    \addlegendentry{$f_1(x)=tan(x)$}

% добавим немного пустого места между описанием первой и второй функций
\addlegendimage{empty legend}\addlegendentry{}

% вторая функция
% здесь необходимо дополнительно ограничить диапазон значений переменной x
    \addplot[blue,samples=256,thick,name path=B] {1-tan(x * 180 / 3.14159)};
    \addlegendentry{$f_2(x)=1-tan(x)$}

% дополнительное пустое место не требуется, так как формулы имеют небольшой размер по высоте

% третья функция
    \addplot[red,samples=256,thick,name path=C] {tan(x * 180 / 3.14159)/10};
    \addlegendentry{$f_3(x)=\frac{tan(x)}{10}$}

% закрашиваем фигуру
\addplot[blue!20,samples=256] fill between[of=A and C,soft clip={domain=0.0:0.4636}];
\addplot[blue!20,samples=256] fill between[of=B and C,soft clip={domain=0.4636:0.7378}];
\addlegendentry{$S=0.166$}

% Поскольку автоматическое вычисление точек пересечения кривых в TiKZ реализовать сложно,
% будем явно задавать координаты.
\addplot[color=black] coordinates {(0.0, 0)} node [label={-135:{\small 0.0}}]{};

\addplot[dashed] coordinates { (0.4636, 0.5) (0.4636, 0) };
\addplot[color=black] coordinates {(0.4636, 0)} node [label={-90:{\small 0.4636}}]{};

\addplot[dashed] coordinates { (0.7378, 0.091) (0.7378, 0) };
\addplot[color=black] coordinates {(0.7378, 0)} node [label={-45:{\small 0.7378}}]{};

\end{axis}
\end{tikzpicture}
\caption{Плоская фигура, ограниченная графиками заданных уравнений}
\label{plot2}
\end{figure}















\newpage
Также, помимо аналитически решенных систем, были протестированы те, которые аналитически не решаются. Например:
\begin{minted}[frame=single, linenos, breaklines=true, fontsize=\small]{text}
0.0 4.0
2 x 4 / tan -
x
0.6 pi *
\end{minted}
и
\begin{minted}[frame=single, linenos, breaklines=true, fontsize=\small]{text}
0.0 4.0
2 x 4 / tan -
x
0.6 pi *
\end{minted}
\begin{table}[h]
\centering
\begin{tabular}{|c|c|c|}
\hline
Кривые & $x$ & $y$ \\
\hline
1 и 2 &  1.5823 & 1.5823 \\
2 и $3_1$ &  2.557 & 1.256 \\
1 и $3_1$ & 1.256 & 1.256 \\
2 и $3_2$ &  0.4581 & 1.885 \\
1 и $3_2$ & 1.885 & 1.885 \\
\hline
\end{tabular}
\caption{Координаты точек пересечения}
\label{table1}
\end{table}

\begin{figure}[h]
\centering
\begin{tikzpicture}
\begin{axis}[% grid=both,                % рисуем координатную сетку (если нужно)
             axis lines=middle,          % рисуем оси координат в привычном для математики месте
             restrict x to domain=-1:3,  % задаем диапазон значений переменной x
             restrict y to domain=-0.25:4.5,  % задаем диапазон значений функции y(x)
             axis equal,                 % требуем соблюдения пропорций по осям x и y
             enlargelimits,              % разрешаем при необходимости увеличивать диапазоны переменных
             legend cell align=left,     % задаем выравнивание в рамке обозначений
             scale=1.9,                    % задаем масштаб 2:1
             xticklabels={,,},           % убираем нумерацию с оси x
             yticklabels={,,}]           % убираем нумерацию с оси y

% первая функция
% параметр samples отвечает за качество прорисовки
    \addplot[green,samples=256,thick,name path=A] {x};
% описание первой функции
    \addlegendentry{$f_1(x)=x$}

% добавим немного пустого места между описанием первой и второй функций
\addlegendimage{empty legend}\addlegendentry{}

% вторая функция
% здесь необходимо дополнительно ограничить диапазон значений переменной x
    \addplot[blue,samples=256,thick,name path=B] {2-tan((x * 180 / 3.14159) / 4)};
    \addlegendentry{$f_2(x)=2-tan(\frac{x}{4})$}

% дополнительное пустое место не требуется, так как формулы имеют небольшой размер по высоте

% третья функция
    \addplot[red,samples=256,thick,name path=C] {0.4 * 3.14159};
    \addlegendentry{$f_{3_1}(x)=0.4 \cdot \pi$}

% добавим немного пустого места между описанием третьей и четвертой функций
\addlegendimage{empty legend}\addlegendentry{}

% четвертая функция
    \addplot[yellow,samples=256,thick,name path=D] {0.6 * 3.14159};
    \addlegendentry{$f_{3_2}(x)=0.6 \cdot \pi$}

% закрашиваем фигуру
\addplot[blue!20,samples=256] fill between[of=A and C,soft clip={domain=1.256:1.5823}];
\addplot[blue!20,samples=256] fill between[of=B and C,soft clip={domain=1.5823:2.557}];
\addlegendentry{$\textcolor{blue}{S_1} = \textcolor{red}{S_2} = 0.166$}

% закрашиваем фигуру
\addplot[red!20,samples=256] fill between[of=B and D,soft clip={domain=0.4581:1.5823}];
\addplot[red!20,samples=256] fill between[of=A and D,soft clip={domain=1.5823:1.885}];

% Поскольку автоматическое вычисление точек пересечения кривых в TiKZ реализовать сложно,
% будем явно задавать координаты.
\addplot[color=black] coordinates {(0.0, 0)} node [label={-135:{\small 0.0}}]{};

\addplot[dashed] coordinates { (0.4581, 1.885) (0.4581, 0) };
\addplot[color=black] coordinates {(0.4581, 0)} node [label={-90:{\small 0.4581}}]{};

\addplot[dashed] coordinates { (1.256, 1.256) (1.256, 0) };
\addplot[color=black] coordinates {(1.256, 0)} node [label={-45:{\small 1.256}}]{};

\addplot[dashed] coordinates { (1.5823, 1.5823) (1.5823, 0) };
\addplot[color=black] coordinates {(1.5823, 0)} node [label={-45:{\small 1.5823}}]{};

\addplot[dashed] coordinates { (2.557, 1.256) (2.557, 0) };
\addplot[color=black] coordinates {(2.557, 0)} node [label={-45:{\small 2.557}}]{};

\end{axis}
\end{tikzpicture}
\caption{Плоская фигура, ограниченная графиками заданных уравнений}
\label{plot2}
\end{figure}

Этот тест показывает, что площадь вычисляется правильно независимо от положения фигуры в пространстве. 


\newpage

\section{Структура программы и спецификация функций}

В данном разделе необходимо привести полный список модулей и функций,
описать их функциональность.

Изобразить графически разбиение программы на компоненты (модули, функции)
и связи между этими компонентами.


\begin{minted}[frame=single, linenos, breaklines=true, fontsize=\small]{text}
src
├── ast.c
├── diff_ast.c
├── main.c
├── numeric.c
└── rpn2asm
    └── rpn2asm.c
\end{minted}


\newpage

\section{Сборка программы (Make-файл)}

В данном разделе необходимо описать зависимости между модулями программы
и привести текст Make-файла. Зависимости проще всего описать диаграммой.

\newpage

\section{Отладка программы, тестирование функций}

В данном разделе необходимо изложить, как именно производилось тестирование
и отладка численных методов. Тестирование предполагает наличие как минимум
трех тестов на каждый из реализованных методов, удовлетворяющих следующим
условиям.
\begin{enumerate}
\item Для данных тестов должно быть возможно аналитически посчитать ответ,
\item Среди тестов должно быть не более одного <<тривиального>> теста
    с точки зрения применяемого метода, то есть такого теста, где порядок
    кривой совпадает с порядком кривой используемой методом для аппроксимации.
\end{enumerate}

Для каждого теста необходимо привести уравнения кривых и нужных производных,
аналитическое вычисление корней и отрезков применения методов, результаты
работы численных методов.

\newpage

\section{Программа на Си и на Ассемблере}

В данном разделе необходимо написать, что
исходные тексты программы имеются в архиве, который приложен к этому отчету.

Наличие приложенного архива с текстами программы является обязательным
требованием сдачи отчета. Если почтовая служба по каким-то причинам блокирует
электронное письмо с вашим архивом, то используйте зашифрованный архив с паролем <<123>>.
Такого пароля достаточно, чтобы почтовые сервисы не заглядывали в архив.

\newpage

\section{Анализ допущенных ошибок}

\newpage
\begin{raggedright}
\addcontentsline{toc}{section}{Список цитируемой литературы}
\begin{thebibliography}{99}
\bibitem{math} Ильин~В.\,А., Садовничий~В.\,А., Сендов~Бл.\,Х. Математический анализ. Т.\,1~---
    Москва: Наука, 1985.
\end{thebibliography}
\end{raggedright}

\newpage

\section*{Требования к оформлению}

В данном разделе приводятся общие требования к оформлению текста отчета.
Данный раздел не должен включаться в сдаваемый отчет.

\begin{enumerate}
\item Отчет оформляется на листах A4. Поля должны составлять от 2 до 4
    сантиметров и быть одинаковыми на всех страницах отчета.
\item Основной текст отчета оформляется пропорциональным шрифтом с засечками,
    таким как Times New Roman. Размер шрифта может составлять либо 12pt, либо 14pt.
    Межстрочные интервалы могут быть единичными или полуторными в случае 12-го шрифта
    и только единичными в случае использования 14-го шрифта.
\item Никаких дополнительных межстрочных интервалов между абзацами не делается.
    Первая строка абзаца должна иметь небольшой отступ (5-10мм), одинаковый для
    всех абзацев, включая первый абзац раздела.
\item Заголовки первого уровня должны быть набраны более крупным шрифтом (16pt или 18pt).
    В заголовках допускается использование как основного шрифта, так и пропорционального
    шрифта без засечек, такого как Arial. Все заголовки всех уровней должны быть набраны
    одним шрифтом. Размер шрифта заголовков большего уровня не должен превосходить размер
    шрифта заголовков меньшего уровня.
\item Фрагменты программ и сценариев сборки должны быть набраны моноширинным шрифтом, таким
    как Courier. Размер шрифта, используемый в листингах программ может отличаться от размера,
    использованного при наборе основного текста, но должен быть одинаковым во всех частях
    отчета и принадлежать интервалу от 10pt до 14pt.
\item Выделение полужирным и/или курсивом допускается для отдельных слов в основном тексте,
    если это требуется. Заголовки рекомендуется выделять жирным.
\item Основной текст выравнивается по двум сторонам. На титульном листе часть текста
    выравнивается по центру, часть по правом краю. Список литературы и названия разделов 
    выравниваются по левому краю.
\item Таблицы и рисунки выравниваются по центру. Все таблицы и рисунки должны быть пронумерованы
    и подписаны. Нумерация сквозная, отдельная для рисунков и таблиц, арабскими цифрами.
\item При использовании растровых изображения для иллюстраций в отчете
    необходимо обеспечить достаточное разрешение этих изображений. Качество изображения
    считается достаточным, если все надписи на нем легко читаются. Если на тексте, содержащемся
    на рисунке, явно заметно размазывание элементов букв, то такое изображение считается
    слишком низкого качества, и оно не должно быть использовано в отчете.
\item Таблицы должны быть сверстаны как таблицы, а не вставлены как рисунки.
\item Список литературы должен содержать для книг и статей (в соответствующем порядке).
    \begin{itemize}
    \item Фамилии и инициалы (либо полные имена) всех авторов.
    \item Название книги или статьи.
    \item Название журнала и номер тома или выпуска для статей.
    \item Город и год издания.
    \end{itemize}
\item Список литературы для электронных источников должен содержать
    \begin{itemize}
    \item Название страницы.
    \item Полный адрес страницы.
    \item Дата обращения.
    \end{itemize}
\item Ссылки на Википедию и другие электронные ресурсы для оценок численных 
    методов не принимаются. Используйте книги и/или научные статьи в качестве 
    источников данной информации.
\item На все элементы списка литературы должны присутствовать ссылки в тексте отчета. Элементы
    списка литературы должны идти в том порядке, в котором ссылки на них первый
    раз встречаются в тексте.
\item Титульный лист оформляется следующим образом.
    \begin{itemize}
    \item Сверху с выравниванием по центру пишется название ВУЗа и факультета. Данный
        фрагмент пишется заглавными или малыми заглавными буквами.
    \item В центре страницы располагается следующая информация (сверху вниз).
        \begin{itemize}
        \item Наименование работы (<<Отчет по заданию №6>>, без кавычек заглавными
            или малыми заглавными буквами).
        \item Тема работы (<<Сборка многомодульных \dots>>, в кавычках, жирным шрифтом).
        \item Вариант. (Без кавычек жирным шрифтом).
        \end{itemize}
    \item Информация о студенте, выполнившем работу и преподавателе выравнивается
        по правому краю. Данный фрагмент набирается обычным шрифтом.
    \item Внизу страницы с выравниванием по центру обычным или немного уменьшенным
        шрифтом пишется город и год выполнения работы.
    \end{itemize}
\end{enumerate}


\end{document}
